\documentclass[a4paper, serif, 10pt]{moderncv}

%% moderncv
% style and color
\moderncvstyle{banking}
\moderncvcolor{black}

%% page layout/margins
\usepackage[top=2.5cm, bottom=2.5cm, left=1.5cm, right=1.5cm]{geometry}

\nopagenumbers{}

%% Babel config for Norwegian language
% ref:
% - https://latex3.github.io/babel/guides/locale-norwegian.html
\usepackage[norsk]{babel}

%% fontspec
\usepackage{fontspec}

\IfFontExistsTF{Linux Libertine}{
  \setmainfont[Ligatures={TeX, Common}, Numbers=OldStyle]{Linux Libertine}
}{}
\IfFontExistsTF{Linux Libertine O}{
  \setmainfont[Ligatures={TeX, Common}, Numbers=OldStyle]{Linux Libertine O}
}{}

%% CTeX
% CJK support, used to show CJK characters in the resume
%
% - fontset=none: disable builtin fontset but instead set the CJK font manually
% - heading=false: disable ctex heading
% - punct=kaiming: use kaiming punctuations styles for CJK
% - scheme=plain: use plain scheme, do not override \normalsize font size
% - space=auto: space settings for CJK characters
%
% ref:
% - http://ctan.mirrorcatalogs.com/language/chinese/ctex/ctex.pdf
\usepackage[UTF8, heading=false, punct=kaiming, scheme=plain, space=auto]{ctex}

\IfFontExistsTF{Noto Serif CJK SC}{
  \setCJKmainfont{Noto Serif CJK SC}
}{}
\IfFontExistsTF{Noto Sans CJK SC}{
  \setCJKsansfont{Noto Sans CJK SC}
}{}

%% line spacing
\usepackage{setspace}
\setstretch{1.125}

%% URL styling - use normal text instead of monospace
\urlstyle{same}

% Auto-underline all links
% Defer until \begin{document} so hyperref is already loaded
\AtBeginDocument{%
  \let\oldhref\href
  \renewcommand{\href}[2]{\oldhref{#1}{\underline{#2}}}%
}

%% Patch \httplink and \httpslink to handle full URLs
% The original moderncv macros always prepend http:// or https:// to URLs.
% This causes issues when the URL already has a protocol (e.g., https://example.com)
% resulting in malformed URLs like https://https://example.com.
% This patch checks if the URL already contains "://" and uses it directly if so.
% Note: We use \str_set:Nx (x-expansion) to expand macros like \@homepage first.
\ExplSyntaxOn
\str_new:N \g_yamlresume_url_str

\RenewDocumentCommand{\httpslink}{O{}m}{
  \str_set:Nx \g_yamlresume_url_str {#2}
  \str_if_in:NnTF \g_yamlresume_url_str {://}
    { \url{#2} }
    { \url{https://#2} }
}

\RenewDocumentCommand{\httplink}{O{}m}{
  \str_set:Nx \g_yamlresume_url_str {#2}
  \str_if_in:NnTF \g_yamlresume_url_str {://}
    { \url{#2} }
    { \url{http://#2} }
}
\ExplSyntaxOff

\name{Ingrid Eriksen}{}
\title{Dataanalytiker med fokus på innsikt og visualisering}
\phone[mobile]{+47 412 34 567}
\email{ingrid.eriksen@gmail.com}
\homepage{https://ingrid-eriksen.no}

\address{Kirkegata 12, Oslo, Oslo, Norge, 0153}{}{}

\extrainfo{{\small \faLinkedin}\ \href{https://www.linkedin.com/in/ingrid-eriksen}{@ingrid-eriksen} {} {} {} • {} {} {}
{\small \faGithub}\ \href{https://github.com/ingeriksen}{@ingeriksen}}

\begin{document}

\maketitle

\section{Grunnleggende informasjon}

\cvline{}{\begin{itemize}
\item Analytisk og nysgjerrig dataanalytiker med erfaring fra offentlig sektor og konsulentbransjen
\item Sterk kompetanse innen dataanalyse, statistikk og maskinlæring, med verktøy som Python, SQL, R og Power BI
\item God til å formidle komplekse sammenhenger til ulike målgrupper og skape beslutningsgrunnlag
\item Erfaring med datavask, visualisering og utvikling av automatiserte analyser
\end{itemize}}
\section{Utdanning}

\cventry{aug. 2016–juni 2021}
        {Master, Samfunnsøkonomi, Poeng: A / 5,0}
        {Universitetet i Oslo}
        {\url{https://www.uio.no/}}
        {}
        {\begin{itemize}
\item Fordypet meg i kvantitative metoder, statistikk og programmering
\item Skrev masteroppgave om sammenhengen mellom kollektivtilbud og boligpriser i Oslo
\item Erfaring med analyse av store datasett i R, Python og Stata
\end{itemize}
\textbf{Kurs}: Statistikk for samfunnsvitere, Økonometri med Python, Maskinlæring for samfunnsvitere, Datavask og datahåndtering, Avansert kvantitativ metode, Samfunnsvitenskapelig dataanalyse}
\section{Arbeidserfaring}

\cventry{sep. 2022–Nå}
        {Dataanalytiker}
        {Statistisk sentralbyrå}
        {\url{https://www.ssb.no}}
        {}
        {\begin{itemize}
\item Utvikler og vedlikeholder analyser av norsk økonomi og befolkning
\item Bygger dashbord og rapporter for interne og eksterne brukere med Power BI og R Markdown
\item Bearbeider og validerer store datasett fra ulike kilder med Python og SQL
\item Presenterer resultater for ledergruppen og bidrar til datadrevne beslutninger
\end{itemize}
\textbf{Nøkkelord}: SQL, Python, Power BI, Statistisk analyse, Datavask}

\cventry{sep. 2021–aug. 2022}
        {Analysekonsulent}
        {Tietoevry}
        {\url{https://www.tietoevry.com}}
        {}
        {\begin{itemize}
\item Gjennomførte kundeanalyser og ad-hoc-rapporter for privat og offentlig sektor
\item Automatiserte gjentakende analyser med Python, noe som reduserte manuelt arbeid med om lag 40 prosent
\item Samarbeidet med prosjektteam for å avklare kundens behov og levere skreddersydde analyseløsninger
\end{itemize}
\textbf{Nøkkelord}: Python, SQL, Excel, Kunderapporter, Automatisering}

\cventry{aug. 2018–juni 2020}
        {Studentassistent i statistikk}
        {Universitetet i Oslo}
        {\url{https://www.uio.no}}
        {}
        {\begin{itemize}
\item Holdt øvingstimer i statistikk og veiledet studenter i bruk av R
\item Bidro til utvikling av oppgaver og eksamenssett i samarbeid med faglærer
\end{itemize}
\textbf{Nøkkelord}: R, Undervisning, Statistikk}
\section{Språk}

\cvline{Norsk}{Morsmål eller tospråklig nivå \hfill \textbf{Nøkkelord}: Morsmål}
\cvline{Engelsk}{Fullt profesjonelt nivå \hfill \textbf{Nøkkelord}: IELTS 7.5, Akademisk og profesjonell}
\section{Ferdigheter}

\cvline{Dataanalyse}{Ekspert \hfill \textbf{Nøkkelord}: Python, SQL, R, Excel, Stata, Pandas, NumPy}
\cvline{Datavisualisering}{Avansert \hfill \textbf{Nøkkelord}: Power BI, Tableau, R Markdown, ggplot2, Matplotlib}
\cvline{Maskinlæring}{Viderekommen \hfill \textbf{Nøkkelord}: scikit-learn, TensorFlow, Keras, NLP, Klassifisering}
\cvline{Statistiske metoder}{Ekspert \hfill \textbf{Nøkkelord}: Regresjonsanalyse, Hypotesetesting, Tidsserieanalyse, Bayesiansk statistikk, Signifikanstesting}
\section{Utmerkelser}

\cventry{sep. 2021}
        {Norsk statistisk forening}
        {Masteroppgaveprisen}
        {}
        {}
        {Fikk pris for masteroppgaven 'Kollektivtilbud og boligpriser i Oslo' – en empirisk analyse med fokus på kausale sammenhenger.}
\section{Sertifikater}

\cventry{mars 2023}
        {Microsoft}
        {Microsoft Certified: Data Analyst Associate}
        {\url{https://learn.microsoft.com/en-us/certifications/data-analyst-associate/}}
        {}
        {}

\cventry{juni 2021}
        {Google}
        {Google Data Analytics Professional Certificate}
        {\url{https://www.coursera.org/professional-certificates/google-data-analytics}}
        {}
        {}
\section{Publikasjoner}

\cventry{mars 2022}
        {Sammenhengen mellom kollektivtilbud og boligpriser: En kvantitativ analyse av Oslo}
        {Samfunnsøkonomen}
        {\url{https://www.samfunnsokonomen.no/}}
        {}
        {\begin{itemize}
\item Analyserer hvordan økt frekvens i kollektivtransporten påvirker boligprisene i Oslo
\item Bruker faste effekter og instrumentvariabler for å identifisere kausale sammenhenger
\item Funnene viser en signifikant, men moderat, positiv effekt på boligpriser i nærheten av holdeplasser
\end{itemize}}
\section{Prosjekter}

\cventry{jan. 2023–apr. 2023}
        {Analyseprosjekt som kombinerer åpne data fra Ruter og SSB for å utforske bruksmønstre og tilgjengelighet.}
        {Oslos kollektivsystem – åpen data-analyse}
        {\url{https://github.com/ingeriksen/kollektiv-analyse}}
        {}
        {\begin{itemize}
\item Utviklet en interaktiv rapport som visualiserer hvordan kollektivtilbudet varierer mellom bydeler
\item Benyttet Python for datavask og geografisk analyse, og bygget dashbord i Plotly Dash
\item Publisert på GitHub med åpen kildekode for gjenbruk
\end{itemize}
\textbf{Nøkkelord}: Python, Plotly, API, Dataanalyse}

\cventry{mai 2021–aug. 2021}
        {En modell for å predikere valgdeltakelse i norske kommuner basert på demografi og tidligere valg.}
        {Valgresultater – prediksjonsmodell}
        {\url{https://github.com/ingeriksen/valgresultater}}
        {}
        {\begin{itemize}
\item Bygget en regresjonsmodell med scikit-learn basert på data fra 356 kommuner
\item Identifiserte utdanningsnivå og alderssammensetning som viktigste prediktorer
\item Dokumenterte prosessen i en Jupyter Notebook med tilhørende forklaringer
\end{itemize}
\textbf{Nøkkelord}: scikit-learn, Jupyter, Regresjon, Valg}
\section{Interesser}

\cvline{Friluftsliv}{Fjellturer, Ski, Kajakk}
\cvline{Teknologi}{Datavisualisering, Åpne data, Kunstig intelligens}
\section{Frivillig arbeid}

\cventry{jan. 2023–Nå}
        {Frivillig datahjelper}
        {Røde Kors Oslo}
        {\url{https://www.rodekors.no}}
        {}
        {\begin{itemize}
\item Hjelper frivillige med å strukturere og analysere medlemsdata
\item Utvikler en enkel oversikt over aktivitetsnivå og rekruttering i Excel og Power BI
\end{itemize}}

\end{document}